\section{Related work}\label{sec:rel}
% src: L6 (all sections), L5 Section 9 and L5-L1 to L5-L15, B6 Sections 3 and 6 (reading levels), B1 to B5 change lists, C2 Section 4; H1 Sections 1 to 9 and its change list, HC Section 4 (literature verdicts against the primary texts), reading levels in the annote fields of references.bib, decisions 18 to 22

This section places the framework among eighteen lines of work and says, for each, what the earlier work established, what we take from it and what we add or change. The first three are treated here, the third in part; the other fifteen are summarized by cluster in \cref{sec:rel-formal} and surveyed in the supplement (\cref{sup:rel}). A statement cited with the mark Supplement is one of the secondary results, extensions, proofs or numerical checks that the supplement holds. Twelve of the eighteen concern hypotheses H1 to H4, extended where \cref{sec:met,sec:hie,sec:ren,sec:env} add to them, and six concern hypothesis H5. Claims about cited works stay at the level at which we read them: where we know a result only through a secondary account, a restatement in a later work or targeted extraction from a full text we say so, and works that we verified only bibliographically are cited for standard results or for attribution alone. Statements that we found no precedent report searches, not proofs of absence, and \cref{sec:rel-new} lists the precedents still to be checked. The survey extends the literature companion \citep{dennis2026literature} of the report that introduced the minimal model of \cref{sec:min} \citep{dennis2026communicating}.

\subsection{The nested-blanket program and its constructions}\label{sec:rel-nested}
% src: L6-F1 to L6-F5, L6 Section 1.5, L2-P9 (literature sentence), B1-T3, B1-T4, B6-F8, B2 summary, C2 Section 4.2 item 2, decisions 11 and 12; H1-R14, H1-R15, H1-R17, H1-R19, H1 Section 9 (H5a), HC Section 4 (Kirchhoff et al., Section 4, and Palacios et al., Eq. 14, verified), HC V5, V7(b), V7(c), decisions 19, 21 and 22; reading levels per the annote fields of references.bib: maisto2025flock Appendix A.4 and waade2025one Section 2 (opening and 2.1) read in primary text

The picture behind hypothesis H1 is stated as a program by \citet{kirchhoff2018markov}, who argue that Markov-blanketed systems assemble into ensembles with blankets of their own, a blanket of blankets that recurs from organelles to societies and behaves like an order parameter; the recursion is argued from the definition of a thing by its blanket, with no theorem, simulation or test of conditional independence. They state one condition: an ensemble blanket can be spoken of only when properties of the collective dynamics feed back into the scale below. \citet{ramstead2018answering} propose, without deriving it, that cultural ensembles minimize free energy within the same nested picture. \citet{palacios2020markov} give the only simulation we found in which a second-level blanket is constructed and coupled to the first: cells with identical priors over the signals of each role infer their roles and form a blanket, and a second level is simulated with one ensemble explicit and fifteen in mean field. The authors call the result a numerical proof of principle that needs the right priors, and conditional independence is read off the limited range of the signals rather than computed. The two levels are coupled by an explicit choice, the one construction of a collective's coarse sensory and active states that we found: the simulated ensemble senses the average of the other ensembles' signals, and the average of its members' active states drives them. \citet{kaufmann2021active} impose the shared goal in their methods although their abstract says that alignment emerges, \citet{heins2023spin,heins2024collective} obtain collective behavior without communication channels or a blanket analysis, \citet{waade2025one} compose agents in a fixed blanket structure, and \citet{maisto2025flock} detect a flock-level blanket after the fact by thresholding the Fiedler vector of a graph that counts shared headings, which assigns whole birds to internal, external, sensory and active roles; that identification computes no conditional independence.

We take the picture of collectives as blankets of blankets and the grading of meta-agenthood, and we change the blanket and add a selection criterion. Roles belong to channels (\cref{def:col-partition}), and under local message passing with private noise and a conditionally independent start every region with a two-hop shell is an exact path blanket (Supplement, \cref{prop:col-two-hop}), so the existence of a blanket selects nothing (Supplement, \cref{prop:col-selects-nothing}). Collectives are selected by contrast, computed on the precision of the stationary world-conditional law, at a level declared in advance above the homogeneous baseline (\cref{def:col-contrast,def:col-baseline,def:col-collective}), and graded by closure, interaction necessity against a declared null family (\cref{def:agy-nulls}) and agency (\cref{thm:agy-h1}). Contrast is a spectral criterion in the line of nearly decomposable systems \citep{simon1961aggregation} and spectral clustering \citep{kannan2004clusterings}. Within active inference, \citet{maisto2025flock} identify a flock-level blanket spectrally from a co-heading graph; our criterion acts instead on the precision of the stationary world-conditional law, at a level declared in advance, with roles assigned to channels. The partitioning theorem of \citet{oveisgharan2014partitioning}, which we know from its abstract, concerns the normalized combinatorial Laplacian: with its eigenvalues ordered as $\varpi_1\le\varpi_2\le\dots$, a gap $\varpi_{J+1}\ge\mathrm{poly}(J)\varpi_J^{1/4}$ lets a spectral algorithm find $J$ sets whose internal conductance is much larger than their conductance to the rest of the graph. Whether a converse of this kind holds for precisions with grounding leaks is open.

Formation is where these simulations and our results meet: in our two-group model (\cref{def:frm-two-group}) roles are not supplied as priors, but the split follows a supplied difference of frames, and how far the result is established is stated with \cref{hyp:int-h3} and in \cref{sec:frm}, and is set beside these simulations in the supplement (\cref{sup:rel-formation-status}).

Recursive particular partitions (Supplement, \cref{sec:rel-renormalization}) and the cellular collectives of \citet{fields2022competency} and \citet{levin2022technological} (the second read through targeted extraction), which act in anatomical morphospace and other problem spaces, also treat a collective as an agent with external states of its own, without conditions under which its states have the properties of an agent. \Cref{sec:met} declares a meta-agent's coarse internal, sensory, active and external states as linear functions of the joint process and gives sufficient conditions, in linear-Gaussian models, under which its sensory and active states form a path blanket (necessary and sufficient ones for the strong coarse path blanket under conditions (b1) to (b3) of \cref{prop:met-sufficiency}(b)), its exterior is exactly closed (by kernel invariance, which is also necessary and for block sums is lumpability) or closed of a stated order, its internal state parametrizes its predictive law of the current exterior under consensus on information contributions, and its joint act follows a group posterior; error bounds exist for all of these, for the predictive law only a conditional one and a finite-horizon one that accumulates over the horizon (Supplement, \cref{thm:apx-posterior,prop:apx-posterior-finite}). The properties hold for every set that meets their hypotheses, so they select nothing: selection stays with contrast and the null family, and hierarchies of collectives are judged level by level (\cref{def:col-hierarchy}).

\subsection{The Markov-blanket debate}\label{sec:rel-debate}
% src: L6-F9 to L6-F12, L6-P1, L6-P2(c),(d) (commutator identity and normal drift, re-derived and numerically confirmed here), L6 Section 2.8, B1-T1, B1-T4, B1-X1, B3-P5, decisions 4 and 14; H1-R7, H1-R9, H1-P1, HC Section 4 (Aguilera et al. verified in arXiv v3; the word renormalization does not occur there)

Whether a blanket can bear the weight that the nested program puts on it is disputed. \citet{bruineberg2022emperors} separate Pearl blankets, properties of a chosen statistical model, from Friston blankets, claimed physical boundaries that would need premises the mathematics does not supply. Our blankets are of the first kind, relative to a declared model and to the external states that members sense (\cref{def:col-partition}), and they license no claim of agency. \citet{biehl2021technical} show that sparse coupling of the drift and a stationary blanket imply each other in neither direction and that the original free energy lemma is false as stated, so we declare the blanket notion in each statement (\cref{def:col-blankets}) and test agency separately (\cref{thm:agy-agency}).

For linear diffusions with drift $-I+gA$ and isotropic noise, \citet{aguilera2022particular} report that sparse coupling gives a stationary blanket to first order in $g$ and violates it at second order through blanket states that drive both sides, except in special cases such as symmetric coupling. The second-order term has a closed form. With the noise normalized so that the stationary covariance is $\sigma^2I$ at $g=0$, the Lyapunov equation reads $\Sigma=\sigma^2I+(g/2)(A\Sigma+\Sigma A^\top)$, and iterating it twice and inverting gives
\[
\Pi=\sigma^{-2}\Bigl[I-\tfrac{g}{2}\bigl(A+A^\top\bigr)+\tfrac{g^{2}}{4}\bigl(A^\top A-AA^\top\bigr)\Bigr]+O(g^{3}).
\]
Under sparse coupling with no blanket state reading both sides, the interior--exterior block of $A+A^\top$ vanishes and that of $A^\top A-AA^\top$ is $-\sum_{b\in\sB}A_{\sI b}A_{\sE b}^\top$, a sum over the blanket states that drive both sides, so the violation is $g^2/(4\sigma^2)$ times a block of the commutator $[A^\top,A]$ and vanishes for normal $A$. For normal $A$ the expansion terminates: the drift is then normal, $\Sigma=-2\sigma^2(J+J^\top)^{-1}$ with $J=-I+gA$ solves the Lyapunov equation, and $\Pi=\sigma^{-2}[I-\frac g2(A+A^\top)]$ exactly, so a sparsely coupled normal drift has an exact stationary blanket at every coupling strength that keeps it stable, with a shell one layer thick and a nonzero stationary current allowed. We checked the expansion numerically on a random six-variable instance, where the residual falls by a factor of about eight each time $g$ is halved (Supplement, \cref{chk:col-sparse-expansion}). This gives the symmetric exception an exact form; its discrete-time counterpart is that reversible level-0 networks, and normal ones with isotropic noise, have world-conditional precisions coupling only agents at most two hops apart (\cref{thm:col-diffusive-precision}). Out of equilibrium, direct tests find no blanket in coupled Lorenz systems, where a Hessian-based treatment had certified one \citep{friston2021stochastic,aguilera2023knitting,aguilera2023confederation}, an exchange we know through secondary accounts. Our counterpart is that non-reciprocal attention can make the world-conditional precision non-diffusive (Supplement, \cref{ex:col-non-diffusive}), and contrast is then undefined (Supplement, \cref{open:col-non-diffusive}).

\citet{friston2023path} obtain, for flows that respect sparse coupling and independent fluctuations, conditional independence of internal and external paths given blanket paths, which \citet{friston2023simpler} state with conditioning on initial states. We adopt path blankets and make them exact for linear-Gaussian dynamics: with per-channel noise, and from a start that the blanket screens, they hold when no sensory channel is a collider, reading both sides in one round (\cref{thm:col-path-blanket} and Supplement, \cref{ex:col-collider}), and from other starts with a density they leak only through the initial slice. This is d-separation in the time-unrolled network \citep{pearl1988probabilistic} with an exact leak and a converse. Under the decay and conditioning hypotheses of \cref{thm:col-boundary}(b), a region with finite interior is an approximate stationary blanket of the world-conditional law once its shell is thick on the scale of an explicit bound on the Markov length (\cref{def:col-markov-length} and Supplement, \cref{thm:col-markov-bound}), and no multiple of the correlation length bounds the Markov length (\cref{prop:col-markov-vs-correlation}), so a shell thicker than the correlation length need not make an approximate blanket (Supplement, \cref{rem:col-shell-error}).

Read in primary text, \citet{aguilera2022particular} find, for these weakly coupled linear systems, the blanket condition and the restrictions on solenoidal flow satisfied only in a narrow region of parameters, and argue that replacing the true flow by its blanket-conditioned average, which equates the rate of change of an average with the average of the rate of change, decouples a trajectory from its own history. The replies strengthen the definition of sparse coupling \citep{friston2022very,friston2023simpler} and concede that it does not imply a blanket in general \citep{heins2022sparse}, positions on sparse coupling that we know secondarily. The paper does not discuss renormalization, but as we read it the argument applies there at every scale, since the elimination of internal states in recursive particular partitions uses the same blanket-conditioned expectation; in stable linear dynamics what that elimination drops includes a memory kernel (\cref{sec:ren}).

\subsection{Active-inference models of communication}\label{sec:rel-communication}
% src: L6-F6 to L6-F8, L6 Section 1.5, L1 Section 9 context paragraph, L1-D3, L1-P3, L1-P14, L1-P16, L1-P20, L1-P24

Earlier active-inference models of communication, which we know through their abstracts and secondary accounts, include the duet of agents with a shared model \citep{friston2015duet,friston2015hermeneutics}, partner choice by expected free energy \citep{albarracin2022epistemic,kastel2023small}, federated inference in an assumed common frame \citep{friston2024federated} and belief sharing read as observation, which produces echo chambers \citep{catal2025belief}; they are discussed in the supplement (\cref{sup:rel-communication-models}).

The process theory is standard \citep{dacosta2020active}. What we add is a model of the receiver: interpretation levels and calibration classes (\cref{def:fw-levels,def:fw-calibration}); reception as prior completion, which is Jeffrey conditioning on the sent coordinates \citep{jeffrey1965logic} and the Kalman gain \citep{kalman1960new} in the Gaussian case (\cref{prop:fw-prior-completion}); and codes identified from the co-occurrence of messages with the receiver's own observations of the world (\cref{prop:fw-grounded-identification}), with symbol grounding \citep{harnad1990symbol} and joint attention \citep{tomasello2005understanding} as conceptual antecedents. Strategic content \citep{crawford1982strategic} is outside our scope.

\subsection{Formal antecedents}\label{sec:rel-formal}

The other fifteen lines of work are summarized here in five clusters, and each paragraph ends with the subsections of the supplement that survey its cluster. Hypothesis H4 rests on bounds for the speed of influence and the decay of correlations. Our light cone (\cref{thm:sca-lightcone}) is the discrete finite-range counterpart of the finite group velocity proved by \citet{lieb1972finite}, and our clustering theorem (\cref{thm:sca-clustering}) is the classical linear analog of exponential clustering \citep{nachtergaele2006lieb,hastings2006spectral}; both are known forms, and what we change is the hypotheses (\cref{asm:sca-standing}). On $\Z^D$ the covariance decays at least at the Ornstein--Zernike rate \citep{ornstein1914accidental}. Fidelity is governed by strong data-processing inequalities \citep{ahlswede1976spreading}: our contraction coefficient is known in substance \citep{makur2020comparison}, and for Gaussian broadcasting on $b$-ary trees we prove a necessary and sufficient condition, where the Kesten--Stigum condition \citep{kesten1966additional} is sufficient in general (Supplement, \cref{sec:rel-lieb,sec:rel-sdpi}).

Hypothesis H2 has precise precedents in estimation. Loopy Gaussian belief propagation \citep{malioutov2006walk,weiss2001correctness} and track fusion, where double counting of common information is classical \citep{barshalom1986effect,julier1997nondivergent,grime1994data}, make the ladder from literal reading to exact fusion known; we place it inside each receiver's generative model (\cref{def:fw-levels,def:fw-calibration}). Consistency is a synchronization problem \citep{singer2011angular,bandeira2013cheeger,abbe2018group}: our holonomy criterion (\cref{thm:coh-holonomy}) imports its form from the author's gauge-theoretic manuscript \citep{dennis2026magent}, and what the framework adds is the origin of the maps, translations learned from grounded co-occurrence (Supplement, \cref{sec:rel-fusion,sec:rel-sync}).

Closure asks when a collective's coarse variables have dynamics of their own. Our elimination of a collective's interior (\cref{thm:agy-mori-zwanzig,thm:agy-closure-identity}) is linear Mori--Zwanzig \citep{zwanzig1961memory,mori1965transport} in discrete time; what we add is a declared closure property (\cref{def:agy-closure}) and its link to contrast (\cref{thm:agy-closure-average}). For stationary linear-Gaussian collectives with full-rank noise and no input, dynamical independence \citep{barnett2023dynamical} and the one-step informational closure of \citet{rosas2024software} coincide in exact form with ours (Supplement, \cref{prop:met-closure-notions}), and near decomposability \citep{simon1961aggregation,simon1962architecture,courtois1977decomposability} is the picture behind the nesting of spectral contrast through a hierarchy of timescales (Supplement, \cref{sec:rel-closure}).

Potential games, opinion dynamics, disorder and the social sciences supply the fourth cluster. Where \citet{wolpert2001optimal} design members' payoffs to serve a world utility, our agency theorem (\cref{thm:agy-agency}) finds, under its protocol with shared belief, an exact potential \citep{monderer1996potential}. On our ring the physical correlation length is the Friedkin--Johnsen influence length \citep{friedkin1990social} (\cref{thm:sca-oz}), which we do not present as new, and social learning names the overvaluation case of H2 several times \citep{demarzo2003persuasion,eyster2010naive,enke2019correlation}. The one-dimensional domain law for conventions, which confirms the Imry--Ma scaling \citep{imry1975random} in one dimension, is a theorem on rings (\cref{thm:coh-domains-1d}), and its higher-dimensional forms are conjectures. The communication tests of \citet{deutsch1953nationalism}, transaction flows \citep{savage1960statistical} and modularity \citep{newman2006modularity} are the closest precedents for contrast and declared nulls (Supplement, \cref{sec:rel-games,sec:rel-opinion,sec:rel-disorder,sec:rel-social}).

For the extension behind hypothesis H5, hierarchical active inference \citep{friston2008hierarchical,friston2017graphical} does not, as we read it, double count, and we found no treatment, there or in the work on cultural priors \citep{veissiere2020thinking,constant2019regimes}, of whether a top-down prior built from the receivers' own contributions counts them twice, which \cref{sec:hie} analyzes in linear-Gaussian models. The renormalization of \citet{friston2019free} lacks, in the parts we read, an account of what the reduction drops, which \cref{sec:ren} supplies for linear-Gaussian models, where Laplacian renormalization \citep{villegas2023laplacian} is shown not to see a uniform precision leak; which regions with a large ratio of internal spectral gap to cut per member that procedure returns is an open problem, a threshold of a constant times the logarithm of the number of agents being excluded by a counterexample (Supplement, \cref{prop:ren-lrg,conj:ren-lrg,prop:obj-lrg-counter}). The selective-sensing results of \cref{sec:env} are anticipated in substance by work on the information bottleneck and rational inattention \citep{chechik2005information,sims2003implications}, and the economics of organizations as information processors \citep{radner1962team,marschak1972economic} is cited for attribution (Supplement, \cref{sec:rel-hierarchy,sec:rel-culture,sec:rel-renormalization,sec:rel-network,sec:rel-sensing,sec:rel-organizations}).

\subsection{What is new, and what remains to be checked}\label{sec:rel-new}
% src: L6 Section 4, L6-O1 to L6-O8, B6-O2, B6-O7, B4-O2; the preprints.org manuscript 202410.1895 named in L6 Section 4 is the preprint of waade2025one (Crossref record of doi 10.20944/preprints202410.1895.v1); H1 Section 9 (new as far as found; to check), H1-O1 to H1-O6, HC Section 6 (open items)

Relative to the active-inference literature we found, at the reading levels recorded in this section and the supplement, the framework is new in combination: collectives selected by contrast at a declared level above a baseline and graded by closure, interaction necessity and agency; receivers modeled at explicit interpretation levels and calibration classes, with H2 as misreading of joint structure; coherence split into fidelity, calibration and consistency, with formation proved only in part, about the frozen-calibration mean-field map of one model and short of \cref{def:frm-target}, as qualified in \cref{hyp:int-h3}; and scale bounds under a moment bound and a named contraction. The extension behind hypothesis H5 adds, in linear-Gaussian models, exact conditions for the properties of a meta-agent, including a predictive law of its coarse exterior, for which we found no precedent, the closure literature being concerned with prediction of the macroscopic process; the double counting of circular hierarchies; error bounds for coarse-graining, with an exact account of what the linear counterpart of the active-inference reduction drops; and sensing under a cap on channels and a budget at a horizon beyond one step, the joint treatment of cap and budget, the budget form and the one-step selection of slow features having precedents (Supplement, \cref{sec:rel-sensing}). The individual mathematical results mostly translate known mathematics into the framework, as the subsections above and the supplement (\cref{sup:rel}) record, and their novelty is limited accordingly. The further extensions developed in the supplement, which \cref{sec:opn} sets against ten open problems posed in the literature, close none of those problems in the generality in which it was posed. One of them qualifies the grading by interaction necessity: in its linear-Gaussian model, when the shared drive is not recorded the members' stationary covariance identifies nothing about interaction, and channel interventions that meet stated invariance conditions can certify it (Supplement, \cref{thm:idn-covariance,thm:idn-intervention}), while an estimator of interaction without a recorded shared drive remains open (\cref{sec:opn-1}). The following precedents remain to be checked.
\begin{itemize}
\item The Gaussian sharpness of the Kesten--Stigum threshold (Supplement, \cref{thm:sca-trees}), which we expect to be known but have not found stated in the reconstruction literature.
\item Exact stationary blankets for normal and reversible dynamics, in continuous time (\cref{sec:rel-debate}) and in discrete time (\cref{thm:col-diffusive-precision}), in the exchange on the free energy principle or in graphical models of stationary Gaussian processes \citep{dahlhaus2000graphical}, which we know through the restatement of \citet{eichler2006fitting}.
\item The reduction of consensus to the holonomies of fundamental cycles (\cref{thm:coh-holonomy}), which we expect to be classical for local systems on graphs.
\item Conditional mutual information or Markov lengths computed for models of opinion dynamics or social learning, and the length $\ell\sqrt{a/\varepsilon}$ in work on consensus with stubborn agents.
\item Whether the aggregation result of \citet{mossel2016efficient}, of which we read the abstract, settles inferential closure under exact fusion, and whether the conditions for group agency of \citet{list2011group}, which we have not read, match our agency property.
\item Primary texts that we know only from records, abstracts or secondary accounts, such as \citet{imry1975random} and, among the works cited in the supplement (\cref{sup:rel}), \citet{vonneumann1950functional}, \citet{halperin1962product}, \citet{degroot1974reaching} and \citet{marchioro1978velocity}, and results known only through restatements, such as Theorem~2.8 of \citet{monderer1996potential}, Theorem~2.2 of \citet{ney1966martin}, Theorem~5.2.1 of \citet{bakry2014analysis}, \citet{lyons1990random} and \citet{kayalar1988error}; these should be read before their results are cited with theorem numbers.
\item Inferential closure of coarse-grained collectives in works we have not read in primary text, \citet{hesp2019multi}, \citet{ramstead2021multiscale} and the appendix of \citet{rosas2024software}; the methods of \citet{friston2021parcels}, in particular its threshold for retaining modes and whether noise correlations between particles are kept; error bounds for adiabatic reductions beyond the papers we searched, starting with \citet{zaghw2026renormalising} and \citet{devries2026active}; and whether the active-inference model of social reality construction of \citet{nomura2026social}, which we have not read, anticipates any part of the combination above.
\item Whether the ascending messages of deep temporal models \citep{friston2017deep} are computed with or without the descending prior; the convergence of iterated learning to the prior \citep{griffiths2007language}; the overweighting of public signals under coordination motives in the text of \citet{morris2002social}; and whether the top-down feedback of \citet{dennis2026formation} removes the child's own contribution.
\item The linear solution of the past-future bottleneck \citep{creutzig2009past}, the texts of \citet{akaike1975markovian}, \citet{chalk2018toward}, \citet{mackowiak2018dynamic} and \citet{lenggenhager2020optimal}, and the scope of \citet{gordon2021relevance}, before a sensing theorem of \cref{sec:env} is called new.
\item Gaussian coarse-graining of interacting units in \citet{meshulam2019coarse} and \citet{bradde2017pca}; leaks and effective resistances under Kron reduction in the text of \citet{dorfler2013kron}; the bounds of near-complete decomposability in \citet{courtois1977decomposability}; and whether the expansions of \citet{geigenmuller1983systematic} give the first-order rate observed numerically for the memoryless reduction (Supplement, \cref{obs:ren-adiabatic-rate}).
\end{itemize}
